% GENERATED FILE. Edit canonical lessons, then run npm run generate:books.
% canonical-source-hash: fnv1a64:0e8277de7f71e397
% canonical-lessons: MR-C349-samasya, MR-C349-adcan, MR-C349-upay, MR-C349-paryay, MR-C349-ukalne, MR-R349-review-342-349

\chapter{Problems and Solutions}
\label{ch:mr-problems-and-solutions}
\hlchaptermodality{349}

\begin{chapteropening}
\textbf{By the end of this chapter:} \emph{I can describe a problem and suggest a way out.}

Complete the last lesson of chapter 349: I can describe a problem and suggest a way out.
\end{chapteropening}

\section[\texorpdfstring{\mr{समस्या} (samasyā)}{समस्या (samasyā)}]{\mr{समस्या} (samasyā) — a problem}

\label{lesson:MR-C349-samasya}



\begin{quote}
Before the new one: say the Marathi for linked, then the Marathi for on the contrary.
\end{quote}

\subsection*{You'll want to know: \mr{समस्या}}
\textbf{\mr{समस्या}} — \emph{samasyā} — \textquotedblleft{}a problem\textquotedblright{}.

\begin{quote}
\mr{शहरात} \mr{पाण्याची} \mr{समस्या} \mr{गंभीर} \mr{आहे}.
\end{quote}

\begin{quote}
\emph{śahrāt pāṇyācī samasyā gambhīr āhe.}
\end{quote}

\begin{quote}
The water problem in the city is serious.
\end{quote}

Feminine.

Formal; in speech \mr{प्रश्न} also means an issue.

\begin{cousinweb}
\textbf{\mr{समस्या}} (\emph{samasyā}) is a tatsama, taken from Sanskrit \emph{samasyā}, \textquotedblleft{}joining together; a half-verse set to be completed\textquotedblright{} (Monier-Williams). Marathi uses it for any problem to be solved.
\end{cousinweb}

\begin{grammarlens}[title={What you've built}]
One more word for giving reasons and opinions.
\end{grammarlens}

\subsection*{Your turn}
\begin{itemize}
\raggedright
  \item \emph{Say it:} \emph{samasyā}
  \item \emph{Say it:} \emph{samasyā}, once more
  \item \emph{Say it:} \emph{śahrāt pāṇyācī samasyā gambhīr āhe.} — The water problem in the city is serious.
  \item \emph{Recall:} say the Marathi for linked, then the Marathi for on the contrary, then say \emph{samasyā} again
\end{itemize}

\begin{tcolorbox}[breakable,colback=teal!4,colframe=teal!35!black,title={Before you move on}]
What does \mr{समस्या} mean? (\textquotedblleft{}A problem\textquotedblright{}.) Say it once more.
\end{tcolorbox}
\section[\texorpdfstring{\mr{अडचण} (aḍcaṇ)}{अडचण (aḍcaṇ)}]{\mr{अडचण} (aḍcaṇ) — a difficulty; a snag}

\label{lesson:MR-C349-adcan}



\begin{quote}
Before the new one: say the Marathi for on the contrary, then the Marathi for a problem.
\end{quote}

\subsection*{You'll want to know: \mr{अडचण}}
\textbf{\mr{अडचण}} — \emph{aḍcaṇ} — \textquotedblleft{}a difficulty; a snag\textquotedblright{}.

\begin{quote}
\mr{नव्या} \mr{शहरात} \mr{सुरुवातीला} \mr{खूप} \mr{अडचणी} \mr{आल्या}.
\end{quote}

\begin{quote}
\emph{navyā śahrāt suruvātīlā khūp aḍcaṇī ālyā.}
\end{quote}

\begin{quote}
There were many difficulties at first in the new city.
\end{quote}

Feminine. Plural: \mr{अडचणी}, \emph{aḍcaṇī}.

\mr{अडचण} \mr{येणे}, a difficulty comes up; \mr{अडचणीत} \mr{असणे}, to be in trouble.

\begin{cousinweb}
\textbf{\mr{अडचण}} (\emph{aḍcaṇ}) is usually traced to the root \emph{*aḍḍ-}, \textquotedblleft{}obstruct, stop\textquotedblright{}, under which CDIAL lists Marathi \emph{aḍṇẽ}, \textquotedblleft{}to be stopped\textquotedblright{} (CDIAL 188).
\end{cousinweb}

\begin{grammarlens}[title={What you've built}]
One more word for giving reasons and opinions.
\end{grammarlens}

\subsection*{Your turn}
\begin{itemize}
\raggedright
  \item \emph{Say it:} \emph{aḍcaṇ}
  \item \emph{Say it:} \emph{aḍcaṇ}, once more
  \item \emph{Say it:} \emph{navyā śahrāt suruvātīlā khūp aḍcaṇī ālyā.} — There were many difficulties at first in the new city.
  \item \emph{Recall:} say the Marathi for on the contrary, then the Marathi for a problem, then say \emph{aḍcaṇ} again
\end{itemize}

\begin{tcolorbox}[breakable,colback=teal!4,colframe=teal!35!black,title={Before you move on}]
What does \mr{अडचण} mean? (\textquotedblleft{}A difficulty; a snag\textquotedblright{}.) Say it once more.
\end{tcolorbox}
\section[\texorpdfstring{\mr{उपाय} (upāy)}{उपाय (upāy)}]{\mr{उपाय} (upāy) — a remedy; a solution}

\label{lesson:MR-C349-upay}



\begin{quote}
Before the new one: say the Marathi for a problem, then the Marathi for a difficulty.
\end{quote}

\subsection*{You'll want to know: \mr{उपाय}}
\textbf{\mr{उपाय}} — \emph{upāy} — \textquotedblleft{}a remedy; a solution\textquotedblright{}.

\begin{quote}
\mr{या} \mr{प्रश्नावर} \mr{एकच} \mr{उपाय} \mr{आहे}.
\end{quote}

\begin{quote}
\emph{yā praśnāvar ekac upāy āhe.}
\end{quote}

\begin{quote}
There is only one remedy for this problem.
\end{quote}

Masculine.

X \mr{वर} \mr{उपाय}, a remedy for X; \mr{उपाय} \mr{करणे}, to take measures.

\begin{cousinweb}
\textbf{\mr{उपाय}} (\emph{upāy}) is a tatsama, taken from Sanskrit \emph{upāya}, \textquotedblleft{}approach, means, stratagem\textquotedblright{} (CDIAL 2308).
\end{cousinweb}

\begin{grammarlens}[title={What you've built}]
One more word for giving reasons and opinions.
\end{grammarlens}

\subsection*{Your turn}
\begin{itemize}
\raggedright
  \item \emph{Say it:} \emph{upāy}
  \item \emph{Say it:} \emph{upāy}, once more
  \item \emph{Say it:} \emph{yā praśnāvar ekac upāy āhe.} — There is only one remedy for this problem.
  \item \emph{Recall:} say the Marathi for a problem, then the Marathi for a difficulty, then say \emph{upāy} again
\end{itemize}

\begin{tcolorbox}[breakable,colback=teal!4,colframe=teal!35!black,title={Before you move on}]
What does \mr{उपाय} mean? (\textquotedblleft{}A remedy; a solution\textquotedblright{}.) Say it once more.
\end{tcolorbox}
\section[\texorpdfstring{\mr{पर्याय} (paryāy)}{पर्याय (paryāy)}]{\mr{पर्याय} (paryāy) — an alternative; an option}

\label{lesson:MR-C349-paryay}



\begin{quote}
Before the new one: say the Marathi for a difficulty, then the Marathi for a remedy.
\end{quote}

\subsection*{You'll want to know: \mr{पर्याय}}
\textbf{\mr{पर्याय}} — \emph{paryāy} — \textquotedblleft{}an alternative; an option\textquotedblright{}.

\begin{quote}
\mr{आमच्याकडे} \mr{दुसरा} \mr{पर्याय} \mr{नव्हता}.
\end{quote}

\begin{quote}
\emph{āmacyākaḍe dusrā paryāy navhatā.}
\end{quote}

\begin{quote}
We had no other option.
\end{quote}

Masculine.

X \mr{ला} \mr{पर्याय} \mr{नाही}, there is no alternative to X.

\begin{cousinweb}
\textbf{\mr{पर्याय}} (\emph{paryāy}) is a tatsama, taken from Sanskrit \emph{paryāya}, \textquotedblleft{}turning round; succession, turn\textquotedblright{} (Monier-Williams). A turn taken in place of another became an alternative.
\end{cousinweb}

\begin{grammarlens}[title={What you've built}]
One more word for giving reasons and opinions.
\end{grammarlens}

\subsection*{Your turn}
\begin{itemize}
\raggedright
  \item \emph{Say it:} \emph{paryāy}
  \item \emph{Say it:} \emph{paryāy}, once more
  \item \emph{Say it:} \emph{āmacyākaḍe dusrā paryāy navhatā.} — We had no other option.
  \item \emph{Recall:} say the Marathi for a difficulty, then the Marathi for a remedy, then say \emph{paryāy} again
\end{itemize}

\begin{tcolorbox}[breakable,colback=teal!4,colframe=teal!35!black,title={Before you move on}]
What does \mr{पर्याय} mean? (\textquotedblleft{}An alternative; an option\textquotedblright{}.) Say it once more.
\end{tcolorbox}
\section[\texorpdfstring{\mr{उकलणे} (ukalṇe)}{उकलणे (ukalṇe)}]{\mr{उकलणे} (ukalṇe) — to unravel; to solve (a puzzle)}

\label{lesson:MR-C349-ukalne}



\begin{quote}
Before the new one: say the Marathi for a remedy, then the Marathi for an alternative.
\end{quote}

\subsection*{You'll want to know: \mr{उकलणे}}
\textbf{\mr{उकलणे}} — \emph{ukalṇe} — \textquotedblleft{}to unravel; to solve (a puzzle)\textquotedblright{}.

\begin{quote}
\mr{मुलाने} \mr{ते} \mr{कोडं} \mr{दहा} \mr{मिनिटांत} \mr{उकललं}.
\end{quote}

\begin{quote}
\emph{mulāne te koḍaṁ dahā miniṭānt ukallaṁ.}
\end{quote}

\begin{quote}
The boy solved the puzzle in ten minutes.
\end{quote}

Past, with -\mr{ने} (he …): \mr{त्याने} \mr{कोडं} \mr{उकललं}, \emph{tyāne koḍaṁ ukallaṁ}.

Also intransitive: \mr{गाठ} \mr{उकलली}, the knot came undone.

\begin{cousinweb}
\textbf{\mr{उकलणे}} (\emph{ukalṇe}) is inherited: CDIAL lists Marathi \emph{ukalṇẽ}, \textquotedblleft{}to be loosened, be disentangled\textquotedblright{}, under \emph{*utkalati}, \textquotedblleft{}goes out\textquotedblright{}, beside Sanskrit \emph{utkalayati}, \textquotedblleft{}loosens\textquotedblright{} (CDIAL 1716).
\end{cousinweb}

\begin{grammarlens}[title={What you've built}]
One more word for giving reasons and opinions.
\end{grammarlens}

\subsection*{Your turn}
\begin{itemize}
\raggedright
  \item \emph{Say it:} \emph{ukalṇe}
  \item \emph{Say it:} \emph{ukalṇe}, once more
  \item \emph{Say it:} \emph{mulāne te koḍaṁ dahā miniṭānt ukallaṁ.} — The boy solved the puzzle in ten minutes.
  \item \emph{Recall:} say the Marathi for a remedy, then the Marathi for an alternative, then say \emph{ukalṇe} again
\end{itemize}

\begin{tcolorbox}[breakable,colback=teal!4,colframe=teal!35!black,title={Before you move on}]
What does \mr{उकलणे} mean? (\textquotedblleft{}To unravel; to solve (a puzzle)\textquotedblright{}.) Say it once more.
\end{tcolorbox}
\section[(dialogue)]{Looking back on reasons, arguments and proof}

\label{lesson:MR-R349-review-342-349}



\begin{quote}
Say the Marathi for a motive and a purpose.
\end{quote}

\subsection*{Your turn}
Say these aloud:

\begin{itemize}
\raggedright
  \item a motive, a purpose, a consequence, an explanation, to come about
  \item a point at issue, logic, a reasoned argument, to put forward, proof
  \item a disagreement, to agree, to object, an inconsistency, a standpoint
  \item a conclusion, to infer, a basis, a summary, in short
  \item to entertain, to elucidate, certainty, to verify, to reject
  \item to oppose, to support, to justify, neutral, against
  \item a connection, a similarity, a context, linked, on the contrary
  \item a problem, a difficulty, a remedy, an alternative, to unravel
\end{itemize}

\begin{tcolorbox}[breakable,colback=teal!4,colframe=teal!35!black,title={Before you move on}]
A motive? (\textbf{\mr{हेतू}}, \emph{hetū}.) A purpose? (\textbf{\mr{उद्देश}}, \emph{uddeś}.) A consequence? (\textbf{\mr{परिणाम}}, \emph{pariṇām}.) An explanation? (\textbf{\mr{स्पष्टीकरण}}, \emph{spaṣṭīkaraṇ}.) To come about? (\textbf{\mr{निर्माण} \mr{होणे}}, \emph{nirmāṇ hoṇe}.) A point at issue? (\textbf{\mr{मुद्दा}}, \emph{muddā}.) Logic? (\textbf{\mr{तर्क}}, \emph{tarka}.) A reasoned argument? (\textbf{\mr{युक्तिवाद}}, \emph{yuktivād}.) To put forward? (\textbf{\mr{मांडणे}}, \emph{māṇḍṇe}.) Proof? (\textbf{\mr{पुरावा}}, \emph{purāvā}.) A disagreement? (\textbf{\mr{मतभेद}}, \emph{matbhed}.) To agree? (\textbf{\mr{सहमत} \mr{होणे}}, \emph{sahmat hoṇe}.) To object? (\textbf{\mr{आक्षेप} \mr{घेणे}}, \emph{ākṣep gheṇe}.) An inconsistency? (\textbf{\mr{विसंगती}}, \emph{visaṅgtī}.) A standpoint? (\textbf{\mr{दृष्टिकोन}}, \emph{dṛṣṭikon}.) A conclusion? (\textbf{\mr{निष्कर्ष}}, \emph{niṣkarṣa}.) To infer? (\textbf{\mr{अनुमान} \mr{काढणे}}, \emph{anumān kāḍhṇe}.) A basis? (\textbf{\mr{आधार}}, \emph{ādhār}.) A summary? (\textbf{\mr{सारांश}}, \emph{sārāṁś}.) In short? (\textbf{\mr{थोडक्यात}}, \emph{thoḍkyāt}.) To entertain? (\textbf{\mr{बाळगणे}}, \emph{bāḷagṇe}.) To elucidate? (\textbf{\mr{विशद} \mr{करणे}}, \emph{viśad karṇe}.) Certainty? (\textbf{\mr{खात्री}}, \emph{khātrī}.) To verify? (\textbf{\mr{पडताळणे}}, \emph{paḍtāḷṇe}.) To reject? (\textbf{\mr{फेटाळणे}}, \emph{pheṭāḷṇe}.) To oppose? (\textbf{\mr{विरोध} \mr{करणे}}, \emph{virodh karṇe}.) To support? (\textbf{\mr{पाठिंबा} \mr{देणे}}, \emph{pāṭhimbā deṇe}.) To justify? (\textbf{\mr{समर्थन} \mr{करणे}}, \emph{samarthan karṇe}.) Neutral? (\textbf{\mr{तटस्थ}}, \emph{taṭastha}.) Against? (\textbf{\mr{विरुद्ध}}, \emph{viruddha}.) A connection? (\textbf{\mr{संबंध}}, \emph{sambandh}.) A similarity? (\textbf{\mr{साम्य}}, \emph{sāmya}.) A context? (\textbf{\mr{संदर्भ}}, \emph{sandarbha}.) Linked? (\textbf{\mr{निगडित}}, \emph{nigḍit}.) On the contrary? (\textbf{\mr{याउलट}}, \emph{yāulaṭ}.) A problem? (\textbf{\mr{समस्या}}, \emph{samasyā}.) A difficulty? (\textbf{\mr{अडचण}}, \emph{aḍcaṇ}.) A remedy? (\textbf{\mr{उपाय}}, \emph{upāy}.) An alternative? (\textbf{\mr{पर्याय}}, \emph{paryāy}.) To unravel? (\textbf{\mr{उकलणे}}, \emph{ukalṇe}.)
\end{tcolorbox}
